From \parencite{adekoya2022applications_ai_em} :

“Several diverse theories, algorithms, and methods are available, which presents a challenge and a barrier in choosing the right AI technique for the appropriate engineering process or manufacturing process and environments.”

“Notwithstanding the benefits of AI techniques in engineering and manufacturing, Wuest et al. (2016) argue that its diversity poses a challenge to countless manufacturing practitioners and engineers regarding the proper technique to select from among the lots.”

“Therefore, it might hinder the utilisation of AI techniques by an appreciable percentage of experts in this field.”

“Hence, there are no fast rules of superiority among ML algorithms.”

“Thus, their performance may differ from one area of application to the other.”

“Besides, the pertinent literature is disseminated over various journals, conference proceedings, and research communities.”

“Hence, conducting a systematic survey to scrutinise and classify the existing literature is worthwhile.”

“However, most of these existing systematic review studies surveyed the applications of a single AI technique [...] Otherwise, cover only one engineering or manufacturing domain [...].”

“Thus, the progress and advancement of AI techniques (like ML and DL) in various engineering and manufacturing systems and environments have not yet been well addressed in the literature.”


From \parencite{wuest2016ml_manufacturing} :

“One area, which saw fast pace developments in terms of not only promising results but also usability, is machine learning.”

“Promising an answer to many of the old and new challenges of manufacturing, machine learning is widely discussed by researchers and practitioners alike.”

“The application of ML techniques increased over the last two decades due to various factors, e.g. the availability of large amounts of complex data [...] and the increased usability and power of available ML tools.”

“ML is already widely applied in different areas of manufacturing, e.g. optimization, control, and troubleshooting.”

“However, the field is very broad and even confusing which presents a challenge and a barrier hindering wide application.”

“However, the field of machine learning is very diverse and many different algorithms, theories, and methods are available.”

“For many manufacturing practitioners, this represents a barrier regarding the adoption of these powerful tools.”

“A major challenge of increasing importance is the question what ML technique and algorithm to choose.”

“There are many different ML methods, tools, and techniques available, each with distinct advantages and disadvantages.”

“The best fitting algorithm has to be found in testing various ones in a realistic environment.”


From \parencite{plathottam2023review_ai_manufacturing} :

“The predictive and analytic capabilities of AI/ML models can provide useful insights to human operators to facilitate planning, support real-time decision-making, and improve manufacturing efficiency and safety.”

“The increasing diversity and complexity of manufacturing tasks, workflows, and supply chains has increased the number of variables and interdependencies.”

“Implementing AI solutions, including using mature AI/ML technologies remains challenging.”

“the choice of which AI/ML solutions to implement for specific manufacturing problems is not always trivial and involves a certain degree of trade-offs.”

“AI/ML implementation decisions must be suited to each company's unique situation and needs”

“pressure from within the industry to use AI has contributed to companies feeling obligated to pursue AI/ML solutions even when they lack a concrete plan.”


From \parencite{weichert2019review_ml_optimization_production}

“there is tremendous progress and large interest in integrating machine learning and optimization methods on the shop floor in order to improve production processes.”

“machine learning [is used] to save energy, time, and resources, and avoid waste.”

“The interest in the integration of machine learning and optimization algorithms in production processes has increased steadily since 2009 and has been at a consistently high level since 2014.”

“The first step to solving a machine learning problem is to identify the data that is available.”

“The knowledge of the structure of production data is crucial for the right choice of machine learning models.”

“The reviewed literature shows that the production-related applications of machine learning with or without optimization are manifold.”

“there is hardly any correlation between the used data, the amount of data, the machine learning algorithms, the used optimizers, and the respective problem from the production.”

“process complexity, the number of data points per input dimension to the model, and the model complexity are highly related.”

“the connection between the process complexity, the stored data amount, and the model complexity [...] was not respected”

“this connection was not respected, leading to complex models being trained on low amounts of data, risking overfitting and/or a lack of interpretability.”


From \parencite{rai2021ml_manufacturing_industry4} :

“The machine learning (ML) field has deeply impacted the manufacturing industry in the context of the Industry 4.0 paradigm.”

“ML techniques enable the generation of actionable intelligence by processing the collected data to increase manufacturing efficiency without significantly changing the required resources.”

“the ability of ML techniques to provide predictive insights has enabled discerning complex manufacturing patterns and offers a pathway for an intelligent decision support system”

“Advances in Big data and Machine Learning (ML) is changing the traditional manufacturing era into the smart manufacturing era of Industry 4.0”

“ML techniques [...] has the potential to become the main driver in uncovering fine-grained intricate production patterns in smart manufacturing paradigm and offering timely decision support”

“While different ML techniques have been used in a variety of manufacturing applications in the past, many open questions and challenges remain”

“Manufacturing is undertaking a definitive shift as it assimilates and is transformed by the usage of machine learning techniques.”

“The utility of ML in manufacturing applications reflects that ML can be integrated into every phase of product lifecycle, starting from design to disposal of the product in a manufacturing environment.”


From \parencite{chen2023ml_manufacturing_industry4} :

“Machine Learning (ML) is playing a critical role in the digitalization of manufacturing operations towards Industry 4.0.”

“Recently, ML has been applied in several fields of production engineering to solve a variety of tasks with different levels of complexity and performance.”

“However, in spite of the enormous number of ML use cases, there is no guidance or standard for developing ML solutions from ideation to deployment.”

“it remains very difficult for non-experts working in the manufacturing industry to begin developing ML solutions for their specific problems.”

“The first challenging part of application is to formulate the actual problems to be solved.”

“the implementation of ML solution for different tasks in an industrial environment from scratch has not yet been fully discussed.”

“A large number of ML use cases have shown the great potential for addressing complex manufacturing problems”

“there are critical challenges to overcome before the successful application of ML in manufacturing can be realized.”

“To fill the gap between academic research and manufacturing industries, we provided an actionable pipeline for the implementation of ML solutions by production engineers from ideation through to deployment”

“We hope that this can provide support in method selection for decision-makers considering ML solutions.”


From \parencite{sonntag2023pattern_ml_product_development} :

“companies lack the implementation of these approaches… due to inadequate knowledge and experience”

“the initial guess of suitable algorithms requires a certain amount of experience.”

“only a fraction of them actually do so”

“75% of the companies surveyed believe that automation approaches… have no significance for them”

“only 16% plan to implement automation approaches”

“lack of experience and knowledge… to recognize opportunities… and identify suitable implementation options”

“To use ML… one has to identify the most suitable ML algorithm out of a large number”

“the procedure is based on… trial-and-error”

“this is… laborious”

“no solution allows for an informed selection of ML algorithms based on a well-defined problem”

“support non-experts in identifying the optimal algorithm”

“make the knowledge… accessible for non-experts”

“pattern language… enabling problem solving by non-experts”

“patterns… make expert knowledge easily accessible”

\parencite{arena2024conceptual_framework_ml_algorithm_selection_predictive_maintenance}

“the need to choose the right algorithm for a specific task arises as a challenging issue”

“there is no unique ML algorithm that can efficiently handle the analysis of every dataset or all use cases”

“this may frequently lead to results that are poor or too generic for the problem at hand”

“it is crucial to define the applicability domain”

“no method is better than all the others for all possible datasets”

“an important step in developing an ML model is deciding which learning method provides the best results”

“most of the articles focus on the development of the ML model… not considering what conditions lead to selecting one rather than another”

“many of them focus on accuracy… prediction speed… or computational requirements”

“no tool relates algorithmic approaches to the maintenance process”

“theoretical accuracy and the computational speed do not ensure satisfactory results”

“making the algorithm selection a non-trivial task”

“implementing an ML-based PdM strategy presents multiple challenges”